\section{Introduction --- 큰 줄기}
\label{sec:intro}

\begin{itemize}
  \item \intent{요즘 GPT Astra(GPT-6 Astra)로 로봇을 명령 제어하는 프로젝트가 많다.}
  \item \intent{물체를 보고 스킬을 생성해서 잡는 방법들이 많지만 실패할 때가 있다. 실패할 때마다 Astra가 개입해서 진행하는 방식을 원한다. 논문으로 제출할 계획이다.}
  \item \intent{예전 VLA에서 상위에 큰 LLM, 하위에 작은 VLM을 둔 것처럼, 상위에 Astra, 하위에 Jev를 둔다.}
  \item \idea{상위 Astra는 느리지만 계획을 세우고, 하위 Jev는 빠르게 객관식으로 판단함}
  \item \intent{Astra의 문제는 API라서 시간이 너무 오래 걸린다는 것이다.}
  \item \intent{Jev(TypeSafe AI, 새로 나옴)는 객관식 문제를 엄청 잘 푼다.}
  \item \intent{Astra로 조정하는 방식이 VLA보다 생각보다 성능이 잘 나오지만, 결국 이미지를 보고 엔드포인트로 움직이는 방식에 그친다.}
  \item \intent{Jev를 쓰면 엔드포인트를 어디로 얼마만큼 움직일지에 대한 순간적인 판단이 몇십 배 빠르다.}
  \item \intent{jev는 지금 사용이 불가능한상황임을 확인했어 애초에 텍스트 기반의 api 활성화 기반자체가 우리랑 엮기에는 약간은 부족할 수 있겠단 생각을 해서 jev를 대체할 무언가 방법을 찾아서 해볼 필요가 있을 것 같아}
  \item \chg{(09-24, 정본 §44, D24) 하위 결정기 = 로컬 VLM typed 선택기 ``Jev-L'': H200에서 vLLM~\cite{vllm2026}으로 Qwen3-VL~\cite{qwen3vl2025}을 돌리고, 생성 없이 보기 토큰 로그 확률을 보기 집합 안에서 재정규화. 입력 = M1 텍스트 상태 + 머리캠 1장. 구조(느린 Astra + 빠른 typed 결정 모델)와 M4·평가는 그대로, 호출 대상만 바뀜. GPT-6 Luna API는 한국 중앙 1.17 s라 3 Hz엔 느림, Claude API는 logprob 없음, 학습 결정 머리는 ``LLM은 학습 안 함''과 부딪쳐 상한 조건으로만}
  \item \idea{(09-24 사전 시험) Jev-L 4B, 텍스트+머리캠, 동시 4개 p95 0.261 s(워밍 방식을 데이터 본 뒤 고친 판, 반복 0.265·0.292; 첫 판 0.819 s도 같이 보고). 사전 규칙상 가장 작은 ``허용'' 모델 = 4B. 배치 불변 켜면 같은 입력의 답·확률 불변}
  \item \intent{하지만 VLA의 엄청난 단점은 다른 환경에 가면 일반화가 안 된다는 것이다. LLM은 일반화가 된다. 이걸 엄청 부각하고 항상 기억한다.}
  \item \intent{안전은 크게 다루지 않는다.}
  \item 파이프라인: \intent{사용자가 명령을 내리면 Astra가 명령을 잘게 쪼개 어떻게 실행할지 미리 계획한다. 계획에 맞는 스킬들을 불러와서 Astra가 진행한다.} \intent{프레임별로 확인해서 Astra의 계획대로 되고 있는지 판별한다.}
  \item \intent{``마차를 이어서는 기차를 만들 수 없다.'' 처음이 틀어지면 아무리 고쳐도 좋은 게 안 나온다.}
  \item \idea{이미 있는 것: Astra 로봇 제어~\cite{agp2026,robodojo2026}(계층 선행은 JEV-Star만), 스텝마다 LLM이 이산 행동을 고르는 방식~\cite{showharness2026}, 멈추지 않고 겹쳐 부르기~\cite{slowbrain2026}, 합의로 앞 구간 확정하기(A3)}
  \item \chg{(09-24) 빠른 Jev 선택과 느린 Astra 비동기 계획의 조합도 이미 있음(JEV-Star~\cite{jevstar2026}, 스타크래프트 2) $\rightarrow$ 계층은 새로움으로 쓰지 않음. 우리 차이는 로봇 실행 중 실패 판정이 부르는 Astra와 M4}
  \item \chg{우리 자리(M4): 시간차로 겹쳐 부른 블랙박스 typed 결정 호출끼리의 합의에, 실행 뒤 코드가 예상한 상태와 측정한 상태의 비교를 더함. 전제가 깨진 표는 버리고 유지·교체·수리(keep/replace/repair) 중 하나를 고르는 확정 규칙}
  \item \chg{주장 문구: ``VLA'' 대신 ``과제 데이터로 미세조정한 학습 정책(VLA 포함)''이라 쓰고, 항상 ``같은 인식 앞단 위에서''를 붙임. 사용자 빨간 줄은 그대로 둠}
  \item \idea{가장 약한 고리: ``같은 술어 위에서 Jev 결정이 룰보다 낫다''는 통제된 조건의 측정을 찾지 못함, 커뮤니티 관측은 오히려 룰이 앞섬(43/50 대 48/50, 신뢰도 낮음) $\rightarrow$ 첫 실험에서 먼저 잼}
  \item \dirn{(09-25 user-log 74·75·77·82) 논문의 중심 틀이 정해짐: VLA를 VLA처럼(과제 전체 end-to-end 정책으로) 쓰지 않는다. 아래 빨간 줄}
  \item \intent{그 아스트라라는 존재를 상위 오케스트라로 쓰는 이유가 존재해야만해 얘가 자기 혼자서도 로봇조작을 할 수 있는놈이기때문에 아스트라를 너무 별일 아닌데도 활용을 해버리면, 그냥 더 작은 vlm쓰는게 맞음 최대한 엄청 최대한 능력치를 다써야해 아스트라를 쓸 수 밖에 없는 이유를 만들 수 있는 장점발휘를 할 수 있게 세팅을 해야함 좀 찾아서 어케 활용할지를 더 구체화해봐봐}
  \item \intent{그리고 한가지 궁금한게 이게 VLA를 따르기는 하는데 VLA는 일반적으로 VLM모델을 엔드투 엔드학ㄱ습을할때 이 엔드포인트를 어디로 움직여야한다는 모델로서의 기능이 아니라 테스크 전체에 활용하는 vlm이잖아? 근데 우리는 이 vlm의 역할을 단순히 오락실 게임기처럼 위아래 앞뒤와 같은 것중에 선택을 하게 하는것에 활용을 한다는 점임을 강ㅇ조를 해야할것 같고 실제로도 그렇게 활용될 수 있도록 하고 있는거지?}
  \item \intent{근데 이게 확실하지는 않잖아? / VLM을 어떤식으로 물체를 향해 로봇을 조종할때 최고의 성능이 나오게 학습할 수 있는지에 대한  신뢰 자료 찾아서 그것도 정리를 좀 해보자ㅣ}
  \item \chg{(user-log 75, 정본 §58·§82) VLM은 오락실 조이스틱: typed 보기(방향·크기 구간·단계·대상)만 고르고 M4가 확정, 연속 action expert가 고른 보기 안의 동작을 채움. 일반 VLA처럼 VLM을 과제 전체의 연속 행동 생성기로 쓰지 않음. 정직하게: 융합 경로에서 expert가 확정 보기를 따르는지는 \tentative{강제도 측정도 안 됨}(코드 투영은 모듈형 경로에만) $\rightarrow$ E-SR0 진단(학습 없음, 조건 GT/예측/반전/무작위) \tentative{예정}. 조사(research/steering\_representation): 동작 수준 명령 인터페이스는 지지(Steerable~\cite{steerable2026} RSS 2026: 단일 형식 중 동작 명령이 가장 효과적, GR 1.5 동작 문장), 반쪽인 점: ``어디로'' 공간 목표 없음(VLA-OS~\cite{vlaos2025} NeurIPS 2025: 시각 근거 표현을 저수준이 더 잘 따름)}
  \item \intent{내가 기존의 VLA 모델에서 문제라고 봤던점은 지금 내가 하는 행동 의 단위 자체가 모든 동작들마다 일관적일 수 없다는 점이야. 생각을 해보자. 우리가 옷을접을때 그냥 책상 위에 있을때는 잘 접을 거야. 그런데, 책상위에 벽돌 30개가 있어서 옷을 잡으러 갈때 계속 방해해버려 VLA는 이걸 절대 못풀어. 그이유는 VLA는 행동의 연속성을 의미단위로 풀려하지만 그 행동들은 의미를 가지기에는 너무 짧은 구간들만 보고 있는거거든 볼거면 전체의 흐름을 모두 본다음에 그 의미와 행동을 매칭해야하는데 그게 아니잖아. 즉 이행동을1}\textasciitilde\intent{3초간의 이미지만 가지고 내가 지금 뭐하고 있는거게}\textasciitilde\intent{ 하는 거나 다름이 없는 이상한 거야. 그래서 효율적인 학습을 위해선 이 전체적인 학습의 흐름이 뭐다 라는걸 아스트라가 잡아주는 과정이 굉장히 중요하다고 생각하는거야. 이걸 어떻게 유기적으로 활용할 수 있을까? 나랑 지금 브레인 스토밍해보자}
  \item \intent{일단 질문을 하기전에, 내가설이 맞는지 틀린지 보정을 해서 확인을 해줘볼래?}
  \item \intent{내가 방해물을 예시로 들었지만 그게 아니라 그냥 다른환경에 취약하다는 얘기야}
  \item \intent{다시 말하자면 새환경에 엄청 취약하다 얘는 그냥 앞뒤 몇개만 가지고 이때는 이러한 흐름을 보니 이미지에서. 그럼 이때는 이런 행동을 해야해 하는데 이게 전체적인 멀리서 보는게 아니라 근시안적인 것들만을 하나로 묶어 학습하는 방식이라는거지 예를 들어 VLM은 의미론적으로 이미지와 텟ㄱ스트를 연결하지만 VLA는 그 액션이라는 것이 이구간에서는 물병을 잡으러갈때 생긴 연결성이었을 수 있지만 다른 구간에서는 선풍기를 잡으러가는 과정이었을 수도 있다는거야. 즉 엄청 멍청하다고보는거야 나는. 그냥 VLA자체가 묶는 하나로 묶어서 학습하는 과정이 무조건 과적합이 될 수 밖에 없는 시스템이라는거}
  \item \idea{(09-25 가설 보정, research/hypothesis\_short\_window) 보정된 가설: 좁고 파편화된 시연 데이터로 순진하게 end-to-end 미세조정하면 VLA는 과제 의미 대신 시점·배치·초기 자세에 묶인 국소 영상$\rightarrow$행동 지름길을 배우는 경향이 강해, 위치·시점·조합이 바뀌면 크게 무너진다(체제 한정 지지, 근거 강함: Shortcut~\cite{shortcut2025} CoRL 2025, KI~\cite{kinsulate2025} NeurIPS 2025; LIBERO-Plus~\cite{liberoplus2025}는 보조). ``무조건 과적합''은 반박($\pi_{0.5}$·KI·GR 1.5가 같은 목적함수로 일반화). ``짧은 창이 원인''은 근거 없음(창을 늘린 이득 +2 pt 이하, 우리 E-TC도 \tentative{+1--4 pt}). 처방 셋 = 데이터 다양성, VLM 지식 보존, 저수준과 정합된 단계·목표 라벨(학습 때만 줘도 이득 상당 부분 남음, ECoT-Lite~\cite{ecotlite2025} CoRL 2025). Astra 자리 = 실행 중 해설자 아님, 새 과제 계획자 + 라벨·규칙 교사(J1·J6). 우리 데이터로 새 환경 취약성은 아직 미시험}
  \item \chg{(09-25 논문) 서론에 보정된 가설을 근거 강도와 함께 넣음(체제 한정·우리 데이터 미시험 명시), 관련 연구에 ``VLM을 조종기로: 조종 표현'' 문단}
  \item \intent{나는 VLA를 VLA 처럼 쓰고 싶지 않은거야 뭔소리인지 알아? 단점이 있기에 그건 순간 순간 어디로 가야하는지 에ㅐㄴㄴ드포인트 정밀 보정하는 역할을 하는 용도로 쓸거여}
  \item \chg{(user-log 82, 논문 중심 틀) VLA = 순간순간의 말단 정밀 보정기(지금 목표 대비 어느 방향으로 얼마만큼). 과제 이해·흐름·어디로는 위에서: Astra J1 과제 컴파일(정본 §82) + 움직임 인터페이스 A+B(M4 느린 투표자 + 물체 기준 경로, user-log 78, \tentative{시험 준비 중}). 채택한 움직임 줄(§83)은 이 보정기의 ``지금 움직이는 중인가'' 입력. 초록·서론·방법 ``조이스틱 원칙''·개요 그림·결론에 반영, 결과 수치는 모두 주황}
\end{itemize}

\paragraph{컨트리뷰션 후보 (v5)}
\intent{컨트리뷰션은 아직 없고 점차 잡아간다.}
\begin{enumerate}
  \item M4: 겹쳐 부른 Jev 답들이 서로 맞는지(a)와 이미 한 동작이 예상대로 됐는지(b)를 함께 보고 확정 (가장 강한 후보. 멈추지 않고 겹쳐 요청하기는 Slow Brain, 합의로 앞 구간 확정하기는 A3가 먼저 함 $\rightarrow$ 새로움은 시간차 호출 합의와, (b)로 전제를 무효화하고 유지·교체·수리를 고르는 것)
  \item \idea{(09-24) 같은 질문을 여러 번 물어 답을 맞춰 봐도, 보기 이름 때문에 생기는 체계적 오류는 걸러지지 않음(보기 이름만 바꿔도 AUC .94$\rightarrow$.23~\cite{typesafe2026}, 심사 전 프리프린트) $\rightarrow$ (a)만으로는 부족하고 (b)가 필요하다는 동기로 씀}
  \item 평가: 같은 인식·같은 실행기 위에서 결정 층만 바꿔, 과제 데이터로 미세조정한 학습 정책(VLA 포함)과 LLM 결정 층이 깨끗한 장면에서 바뀐 환경으로 갈 때 얼마나 떨어지는지 짝지어 잰 연구는 없음(2026-09-23 제출분 기준). RoboDawn·RoboDojo는 뒷받침하는 선행
  \item M8: 멈추지 않는 로봇에서 호출 시점과 effort를 함께 평가 (평가 쪽 기여. 정체가 길어지면 부르는 방식은 이미 있음) \chg{(09-24) 논문을 한 편으로 쓰므로 첫 논문에서는 최소판. effort는 low가 기본이고 low·high를 비교}
  \item M6: 스킬 안의 결정 지점을 typed 질문으로 만듦 (로봇 쪽 선점 없음, 로봇 밖 선례는 있음)
  \item M10: Astra가 정리한 규칙만 Jev에 넣음 (규칙 정리 자체는 선행이 있어 새로움은 좁음)
\end{enumerate}
\chg{v2의 ``호출 방식 비교가 새롭다''는 틀렸음. 주기 호출과 이벤트 호출을 비교한 연구가 이미 여러 편 있음~\cite{checkvla2026,eventtrig2026}}

\paragraph{논문 틀 (모의 심사 2026-09-23, 사용자 결정 09-24)}
\begin{itemize}
  \item \idea{모의 심사 3명: 3 / 4 / 2점. 가장 큰 지적은 ``두 기여가 서로 다른 시험대에서 따로 검증돼 사실상 두 논문''이라는 것}
  \item \dirn{1안(나눔): (가) CVPR 2027 분석 논문 ``같은 인식 위에서 학습 정책 대 학습 없는 LLM 결정 층'' (나) M4 확정 규칙만 다룬 로봇 논문}
  \item \dirn{2안(한 논문): M4를 방법으로, 일반화를 평가로. 두 기여를 잇는 실험(새 장면 실시간 트랙에서 M4 켬/끔)이 꼭 필요}
  \item \idea{CVPR 2027 제출 마감 2026-11-16(공식 확인). ICRA 2027은 이미 마감}
  \item \intent{논문을 일단 한 편으로 작성을 할 생각이야}
  \item \chg{(09-24 사용자 결정) 한 편 = 2안. 두 기여를 잇는 실험(새 장면 실시간 트랙에서 M4 켬/끔)과 지연 스윕 실험이 필수. M8 호출 방식 비교, M9 복구, M10 경험 축적은 첫 논문에서 최소판}
  \item \decide{실물 로봇 실험 여부}
\end{itemize}
